\section{Variables conjuntas discretas (la tabla)}

\begin{receta}{cualquier problema con tabla}
\begin{enumerate}
\item \textbf{Ubica qui\'en es qui\'en:} qu\'e variable va en renglones y cu\'al en columnas. Escr\'ibelo.
\item \textbf{Agrega los m\'argenes:} total de cada rengl\'on y de cada columna. Verifica que todo sume 1.
\item \textbf{Probabilidad de un evento} ($X>Y$, $X+Y=4$, $X=Y$\ldots): marca las celdas que cumplen y s\'umalas.
\item \textbf{Condicional:} toma el rengl\'on o columna del valor dado y divide cada celda entre su total.
\item \textbf{Independencia:} compara cada celda con (total de su rengl\'on)$\times$(total de su columna).
\end{enumerate}
\end{receta}

\begin{ejemplo}{Ejemplo resuelto 1 --- L\'ineas de producci\'on \origen{Apuntes, Ej.~3.3}}
$X$ = art\'iculos de la l\'inea uno, $Y$ = art\'iculos de la l\'inea dos.
\begin{center}
\begin{tabular}{cc|ccc|c}
 & & \multicolumn{3}{c|}{$x$} & \\
 & $\fxy$ & 0 & 1 & 2 & $f_Y(y)$\\ \hline
 \multirow{3}{*}{$y$} & 0 & 0.10 & 0.20 & 0.20 & \textbf{0.5}\\
 & 1 & 0.04 & 0.08 & 0.08 & \textbf{0.2}\\
 & 2 & 0.06 & 0.12 & 0.12 & \textbf{0.3}\\ \hline
 & $f_X(x)$ & \textbf{0.20} & \textbf{0.40} & \textbf{0.40} & 1
\end{tabular}
\end{center}
\textbf{a) $P(X>Y)$.} Celdas con $x>y$: $(1,0)$, $(2,0)$, $(2,1)$:
\[P(X>Y)=0.20+0.20+0.08=\resp{0.48}\]
\textbf{b) Marginales.} Son los m\'argenes de la tabla (en negritas).

\textbf{c) Condicional de $X$ dado $Y=1$.} Rengl\'on $y=1$ entre su total $0.2$:
\begin{center}
\begin{tabular}{c|ccc}
$x$ & 0 & 1 & 2\\ \hline
$f_{X\mid Y=1}(x)$ & $0.04/0.2=0.2$ & $0.08/0.2=0.4$ & $0.08/0.2=0.4$
\end{tabular}
\end{center}
\textbf{d) $P(X<2\mid Y=1)$} $=0.2+0.4=\resp{0.6}$.

\textbf{e) ?`Independientes?} Cada celda es igual al producto de sus m\'argenes
($0.10=0.2\cdot0.5$, $0.08=0.4\cdot0.2$, \ldots, las nueve). \resp{\text{S\'i son independientes}}.
Observa que por eso la condicional del inciso c) sali\'o igual a la marginal $f_X$.
\end{ejemplo}

\begin{ejemplo}{Ejemplo resuelto 2 --- Supermercado \origen{Serie 3, problema 1}}
$X_1$ = clientes en caja r\'apida (renglones, 0 a 4), $X_2$ = clientes en caja superveloz (columnas, 0 a 3).
\begin{center}\small
\begin{tabular}{c|cccc|c}
$x_1\backslash x_2$ & 0 & 1 & 2 & 3 & $f_{X_1}$\\ \hline
0 & 0.08 & 0.07 & 0.04 & 0.00 & 0.19\\
1 & 0.06 & 0.15 & 0.05 & 0.04 & 0.30\\
2 & 0.05 & 0.04 & 0.10 & 0.06 & 0.25\\
3 & 0.00 & 0.03 & 0.04 & 0.07 & 0.14\\
4 & 0.00 & 0.01 & 0.05 & 0.06 & 0.12\\ \hline
$f_{X_2}$ & 0.19 & 0.30 & 0.28 & 0.23 & 1
\end{tabular}
\end{center}
\textbf{a)} $P(X_1=1,X_2=1)=\resp{0.15}$ (una sola celda).

\textbf{b)} $P(X_1=X_2)$ es la diagonal: $0.08+0.15+0.10+0.07=\resp{0.40}$.

\textbf{c)} $A=\{|X_1-X_2|\ge 2\}$. Celdas: $(0,2),(0,3),(1,3),(2,0),(3,0),(3,1),(4,0),(4,1),(4,2)$:
\[P(A)=0.04+0+0.04+0.05+0+0.03+0+0.01+0.05=\resp{0.22}\]
\textbf{d)} $P(X_1+X_2=4)$: celdas $(1,3),(2,2),(3,1),(4,0)$: $0.04+0.10+0.03+0=\resp{0.17}$.

$P(X_1+X_2\ge4)$: se agregan $(2,3),(3,2),(3,3),(4,1),(4,2),(4,3)$:
$0.17+0.06+0.04+0.07+0.01+0.05+0.06=\resp{0.46}$.
\end{ejemplo}

\begin{trampa}
\begin{itemize}
\item Leer la tabla al rev\'es. En la Serie~1 los renglones son $x_1$; en el Ej.~3.3 los renglones son $y$. Revisa siempre el encabezado.
\item \guillemotleft A lo sumo\guillemotright{} es $\le$, \guillemotleft por lo menos\guillemotright{} es $\ge$, \guillemotleft m\'as de\guillemotright{} es $>$.
\item Para \guillemotleft independientes\guillemotright{} basta \textbf{una} celda que falle para decir que no; para decir que s\'i debes revisar \textbf{todas}.
\end{itemize}
\end{trampa}

\begin{ejemplo}{Ejemplo resuelto 3 --- Restaurante \origen{Serie 3, problema 8}}
$X$ = costo del hombre (renglones), $Y$ = costo de la mujer (columnas).
\begin{center}
\begin{tabular}{c|ccc|c}
$x\backslash y$ & 7 & 9 & 10 & $f_X$\\ \hline
7 & 0.05 & 0.05 & 0.10 & 0.20\\
9 & 0.05 & 0.10 & 0.35 & 0.50\\
10 & 0 & 0.20 & 0.10 & 0.30\\ \hline
$f_Y$ & 0.10 & 0.35 & 0.55 & 1
\end{tabular}
\end{center}
\textbf{a)} Marginales: los m\'argenes.

\textbf{b)} $P(X\le9,\,Y\le9)=0.05+0.05+0.05+0.10=\resp{0.25}$.

\textbf{c)} $\fxy(10,7)=0$ pero $f_X(10)f_Y(7)=0.3\cdot0.1=0.03\neq0$. \resp{\text{No son independientes}}.

\textbf{d)} Costo total esperado: $\E(X+Y)=\E(X)+\E(Y)$.
\begin{align*}
\E(X)&=7(0.2)+9(0.5)+10(0.3)=8.90\\
\E(Y)&=7(0.10)+9(0.35)+10(0.55)=9.35
\end{align*}
$\E(X+Y)=\resp{18.25}$.

\textbf{e)} Devoluci\'on $=|X-Y|$. Se multiplica $|x-y|$ por cada celda (la diagonal aporta 0):
\[
\E|X-Y|=2(0.05)+3(0.10)+2(0.05)+1(0.35)+3(0)+1(0.20)=\resp{1.05}
\]
\end{ejemplo}

\begin{mecaniza}
Con la tabla del restaurante obt\'en la funci\'on condicional de $X$ dado $Y=9$ y $P(X\ge 9\mid Y=9)$.

\textbf{Respuesta:} columna $y=9$ entre $0.35$: $f_{X\mid Y=9}=\{7\!:\tfrac17,\;9\!:\tfrac27,\;10\!:\tfrac47\}$;
$P(X\ge9\mid Y=9)=\tfrac67\approx0.857$.
\end{mecaniza}

\subsection{Funci\'on de distribuci\'on conjunta (tabla acumulada)}

$F_{XY}(x,y)=P(X\le x,\,Y\le y)$. En una tabla: cada celda de $F$ es la suma de todas las celdas de $f$ que quedan
\textbf{arriba y a la izquierda} de ella, incluy\'endola.

\begin{ejemplo}{Ejemplo resuelto 4 \origen{Apuntes, Ej.~3.7 (en blanco en tus apuntes)}}
Con la tabla del Ej.~3.3:
\begin{center}
\begin{tabular}{cc|ccc}
 & $F_{XY}$ & $x=0$ & $x=1$ & $x=2$\\ \hline
 & $y=0$ & 0.10 & 0.30 & 0.50\\
 & $y=1$ & 0.14 & 0.42 & 0.70\\
 & $y=2$ & 0.20 & 0.60 & 1.00
\end{tabular}
\end{center}
Por ejemplo $F(1,1)=0.10+0.20+0.04+0.08=0.42$. Comprobaciones: la esquina final siempre vale 1; el \'ultimo
rengl\'on es $F_X$ y la \'ultima columna es $F_Y$.
\end{ejemplo}

\subsection{Construir la tabla desde el experimento}

\begin{ejemplo}{Ejemplo resuelto 5 --- Transistores \origen{Serie 3, problema 10a}}
Cinco transistores, tres defectuosos (D) y dos buenos (B). $X$ = pruebas hasta el primer D;
$Y$ = pruebas adicionales hasta el segundo D. Se multiplica rama por rama, sin reemplazo:
{\small\begin{align*}
P(X=1,Y=1)&=P(DD)=\tfrac35\cdot\tfrac24=\tfrac{3}{10} &
P(X=1,Y=2)&=P(DBD)=\tfrac35\cdot\tfrac24\cdot\tfrac23=\tfrac{2}{10}\\
P(X=1,Y=3)&=P(DBBD)=\tfrac35\cdot\tfrac24\cdot\tfrac13\cdot1=\tfrac{1}{10} &
P(X=2,Y=1)&=P(BDD)=\tfrac25\cdot\tfrac34\cdot\tfrac23=\tfrac{2}{10}\\
P(X=2,Y=2)&=P(BDBD)=\tfrac25\cdot\tfrac34\cdot\tfrac13\cdot1=\tfrac{1}{10} &
P(X=3,Y=1)&=P(BBDD)=\tfrac25\cdot\tfrac14\cdot1\cdot1=\tfrac{1}{10}
\end{align*}}
\begin{center}
\begin{tabular}{c|ccc|c}
$x\backslash y$ & 1 & 2 & 3 & $f_X$\\ \hline
1 & 0.3 & 0.2 & 0.1 & 0.6\\
2 & 0.2 & 0.1 & 0 & 0.3\\
3 & 0.1 & 0 & 0 & 0.1\\ \hline
$f_Y$ & 0.6 & 0.3 & 0.1 & 1
\end{tabular}
\end{center}
La suma da 1, as\'i que no falta ning\'un caso. Esta tabla se vuelve a usar en la secci\'on de correlaci\'on.
\end{ejemplo}
